% GENERATED by build_chapters.py; edit docs/report/chapters, not this file

\chapter{The n8n Healing Workflow}

The workflow "Self-Healing CI/\allowbreak{}CD Pipeline (Qwen3.8)" has 26 nodes (source: n8n/\allowbreak{}self\_\allowbreak{}healing\_\allowbreak{}pipeline.json). Most of the time it is a straight line of nodes with two forks. This chapter follows one run through seven stages and then explains the retry and re-run behaviour. It does not repeat every node setting. The node-by-node reference, with each node's type, endpoint, error setting and tests, is in \texttt{docs/\allowbreak{}architecture/\allowbreak{}n8n-\allowbreak{}workflow.\allowbreak{}md}. We checked the points below against the workflow JSON, \texttt{tools/\allowbreak{}gen\_\allowbreak{}n8n.\allowbreak{}py} and \texttt{tools/\allowbreak{}nodes/\allowbreak{}*.\allowbreak{}js}.

Two facts help when reading the stages. The JSON is generated by \texttt{tools/\allowbreak{}gen\_\allowbreak{}n8n.\allowbreak{}py} from node specifications, and the Code node bodies are separate files in \texttt{tools/\allowbreak{}nodes/\allowbreak{}}. And a node marked \texttt{onError:\allowbreak{} continueRegularOutput} does not stop the run when it fails. It passes an item with an \texttt{error} key to the next node, which is how failures are routed to the reporting stage.

\section{Stage 1: receive}

The entry point is a POST webhook on the path \texttt{github-\allowbreak{}failure}. It uses header authentication with the credential "N8N Webhook Secret", and its response mode is \texttt{onReceived}, so GitHub gets its answer at once and n8n keeps working in the background (source: n8n/\allowbreak{}self\_\allowbreak{}healing\_\allowbreak{}pipeline.json). The body is the JSON the notify job sends: \texttt{run\_\allowbreak{}id}, \texttt{repo}, \texttt{branch}, \texttt{commit}, \texttt{actor}, \texttt{failure\_\allowbreak{}kind}, \texttt{attempt} and \texttt{pr\_\allowbreak{}number} (Chapter 4 lists them).

There is a second, separate webhook on the path \texttt{set-\allowbreak{}qwen-\allowbreak{}url}. It handles the model address and takes no part in a heal. Stage 4 covers it.

\section{Stage 2: gather evidence}

Six nodes collect what the model will be shown.

\begin{enumerate}[leftmargin=*, itemsep=2pt]
\item \textbf{Fetch Job Logs} lists the jobs of the failed run and their steps. It has no error setting, so if GitHub refuses it the run stops here.
\item \textbf{Fetch Failed Job Log} asks for the log of the first job whose conclusion is \texttt{failure}. It is set not to follow redirects, because GitHub answers with a 302 to a signed address on another host.
\item \textbf{Download Job Log} fetches that signed address with no credential, so the GitHub token is never sent to the storage host (Chapter 4 explains why this was split in two). If the log is missing or expired, the run goes on and the prompt says \texttt{(log unavailable)} (source: tools/\allowbreak{}nodes/\allowbreak{}build\_\allowbreak{}ai\_\allowbreak{}prompt.js).
\item \textbf{Fetch Changed Files} reads the commit's file list.
\item \textbf{Pick Candidate Files} and \textbf{Fetch File Content} choose and download the source.
\end{enumerate}

The choice of files is a small piece of logic in \texttt{tools/\allowbreak{}nodes/\allowbreak{}pick\_\allowbreak{}candidate\_\allowbreak{}files.\allowbreak{}js}. It looks for file paths under the runner's checkout folder in the log, keeps only code files outside \texttt{node\_\allowbreak{}modules}, and adds the code files the commit changed. The order is the first two files from the stack trace, then the changed files, then any other traced files, with duplicates removed and a limit of 3 files (source: tools/\allowbreak{}nodes/\allowbreak{}pick\_\allowbreak{}candidate\_\allowbreak{}files.js, \texttt{MAX = 3}). Using the trace matters because the file that crashes is not always the file that was edited. If nothing qualifies, the node emits one placeholder item so that the run still reaches the failure report, since n8n ends a run quietly when a node outputs nothing.

File contents are fetched at the \texttt{commit} from the webhook. On a retry that is the head commit of the fix PR, so the model sees the file with its own earlier fix already applied.

\section{Stage 3: build the prompt}

\texttt{Build AI Prompt} (tools/\allowbreak{}nodes/\allowbreak{}build\_\allowbreak{}ai\_\allowbreak{}prompt.js) assembles the whole model request in one Code node, so n8n never has to splice a multi-line string into a JSON body.

First it applies the failure gate: if \texttt{failure\_\allowbreak{}kind} is set and is not \texttt{code}, it throws and the run ends with no model call. Then it works out whether this is a retry, which is true only when \texttt{attempt} is above 0 and the branch starts with \texttt{ai-\allowbreak{}fix-\allowbreak{}}.

The CI log is trimmed before it goes in. The node strips timestamps and colour codes, finds the first \texttt{\#\#[error]} line, starts after the closest \texttt{\#\#[endgroup]} above it, and keeps the last 40 lines and at most the last 3000 characters (source: tools/\allowbreak{}nodes/\allowbreak{}build\_\allowbreak{}ai\_\allowbreak{}prompt.js, \texttt{LOG\_\allowbreak{}MAX\_\allowbreak{}LINES} and \texttt{LOG\_\allowbreak{}MAX\_\allowbreak{}CHARS}). Presumably this keeps the prompt short, since the model only needs the part of the log that explains the failure.

The system prompt asks for one raw JSON object with \texttt{branch\_\allowbreak{}name}, \texttt{file\_\allowbreak{}path}, \texttt{edits}, \texttt{pr\_\allowbreak{}title} and \texttt{pr\_\allowbreak{}body}, where \texttt{edits} holds SEARCH/\allowbreak{}REPLACE blocks. The request body uses the model name \texttt{local}, temperature 0.1, \texttt{max\_\allowbreak{}tokens} 4096, \texttt{stream:\allowbreak{} true}, and \texttt{enable\_\allowbreak{}thinking} set to the retry flag. On a retry the prompt also states that the file already contains the previous automated fix and that CI still fails.

\section{Stage 4: call the model, with a fallback}

\textbf{Get Model URL} builds the request address from the value stored in the workflow's static data, and falls back to the variable \texttt{QWEN\_\allowbreak{}BASE\_\allowbreak{}URL} (source: tools/\allowbreak{}nodes/\allowbreak{}get\_\allowbreak{}model\_\allowbreak{}url.js). The stored value comes from the second webhook. After each Kaggle boot the launcher posts the new tunnel address to \texttt{set-\allowbreak{}qwen-\allowbreak{}url}, and \textbf{Store Model URL} saves it only if it matches \texttt{https:\allowbreak{}/\allowbreak{}/\allowbreak{}\textless{}name\textgreater{}.\allowbreak{}trycloudflare.\allowbreak{}com}.

\textbf{Qwen3.8 Analyze} posts to \texttt{/\allowbreak{}v1/\allowbreak{}chat/\allowbreak{}completions} with the endpoint key and a timeout of 240 seconds (source: n8n/\allowbreak{}self\_\allowbreak{}healing\_\allowbreak{}pipeline.json). It reads the response as text, so the stream arrives as one string. The 240 s limit is there so that a stuck endpoint reaches the fallback while the run is still useful (source: comment in tools/\allowbreak{}gen\_\allowbreak{}n8n.py, as reported in docs/\allowbreak{}architecture/\allowbreak{}n8n-workflow.md).

The node \textbf{Model OK?} checks the result. The answer counts as good when there is no error and the text contains \texttt{[DONE]} or a \texttt{finish\_\allowbreak{}reason}. Anything else, such as a refused connection, a timeout or a stream cut short, goes to \textbf{OpenAI Fallback}, which sends the same two messages to \texttt{gpt-\allowbreak{}4.\allowbreak{}1}. \textbf{Normalize OpenAI Reply} converts that answer to the shape the next node expects, or throws one message that names both failures.

Both paths meet at \textbf{Assemble Model Reply}, which rebuilds a streamed answer into the usual completion object. It rejects a stream that ended early, an empty one, and one with an error event.

The fallback has a documented weakness. In the one test on n8n Cloud, using the free OpenAI credits that come with the trial, it answered \texttt{Forbidden} for runs started by the webhook (source: docs/\allowbreak{}guides/\allowbreak{}cloud-n8n.md). So we treat it as untested until a real OpenAI key is attached.

\section{Stage 5: extract the fix}

\texttt{Extract AI Response} (tools/\allowbreak{}nodes/\allowbreak{}extract\_\allowbreak{}ai\_\allowbreak{}response.js) turns the reply into a patched file. It is deliberately strict, because a quiet mistake here becomes a pull request.

\begin{enumerate}[leftmargin=*, itemsep=2pt]
\item It removes any \texttt{\textless{}think\textgreater{}} block and stray code fences, then parses the JSON.
\item It selects the file named in \texttt{file\_\allowbreak{}path}, and only accepts a file that was among those shown to the model.
\item It applies each SEARCH/\allowbreak{}REPLACE block in order. Every SEARCH text must be non-empty and must occur exactly once in the file as patched so far. A block that is not found, matches twice or is malformed rejects the whole answer.
\item It rejects the answer if the result equals the original file.
\item It parses the patched text with \texttt{new Function(.\allowbreak{}.\allowbreak{}.\allowbreak{})}, which checks the syntax without running the code. If the parse fails, no pull request is opened.
\end{enumerate}

The parse check has a known limit, which its comment states plainly: a fix that parses but is wrong still goes through, and the pull request's own CI is what catches it. We could not confirm how the check behaves for a non-JavaScript file, since it uses a JavaScript parser; the demo app is Node, so this never came up.

The node then chooses the branch name. If the model gave none, or \texttt{main} or \texttt{master}, the node uses \texttt{ai-\allowbreak{}fix-\allowbreak{}} plus a timestamp. It appends the first 7 characters of the commit. On a retry it replaces the name with the branch from the webhook. The commit message is the PR title, a blank line and \texttt{Heal-\allowbreak{}Attempt:\allowbreak{} \textless{}attempt + 1\textgreater{}}.

\textbf{Fix Ready?} is the fork. If there is no error and the file content exists, the run continues to the pull request. Otherwise it goes to the failure report (Stage 7).

\section{Stage 6: write the pull request}

Six nodes act on GitHub in this order.

\begin{enumerate}[leftmargin=*, itemsep=2pt]
\item \textbf{Create Branch} posts a new ref from the failing commit. A 422 "Reference already exists" is allowed to pass as data.
\item \textbf{Branch Ready} accepts that case and stops on any other error, such as a bad token or a rate limit, instead of reporting a success it did not reach.
\item \textbf{Get File SHA} reads the file on the heal branch, because updating an existing file needs its blob SHA.
\item \textbf{Push Fix File} writes the patched file with the commit message above.
\item \textbf{Open Pull Request} opens a PR from the heal branch into the branch that failed.
\item \textbf{Heal Outcome} reads the answer and reports \texttt{pr\_\allowbreak{}opened}, \texttt{retry\_\allowbreak{}pushed} or \texttt{pr\_\allowbreak{}already\_\allowbreak{}open}. Anything else throws, so a real failure is not recorded as a success.
\end{enumerate}

The four HTTP nodes among them use the "GitHub PAT" credential. Create Branch and Open Pull Request continue on error. Get File SHA and Push Fix File stop the run when they fail.

\section{Stage 7: report failures}

When no fix is produced, \textbf{Explain Failure} builds a comment and \textbf{Post Failure Comment} posts it. The comment goes on the heal PR if \texttt{pr\_\allowbreak{}number} is all digits, and on the failing commit otherwise. It starts with "Self-healing pipeline: no fix produced" and quotes the reason, cut to 600 characters (source: tools/\allowbreak{}nodes/\allowbreak{}explain\_\allowbreak{}failure.js).

This stage exists because of a failure found in testing. Before it, a retry that died at the tunnel's time limit left a red pull request with no explanation (source: docs/\allowbreak{}results/\allowbreak{}GOLIVE\_\allowbreak{}RESULTS.md, Round 3, "Known limits"). In the later round, the stage was proved with the model endpoint switched off: a comment appeared on the failing commit, and a comment appeared on the PR for a retry (source: docs/\allowbreak{}results/\allowbreak{}GOLIVE\_\allowbreak{}RESULTS.md, Round 5, Fix 2).

The stage has gaps. Several nodes have no error setting (Fetch Job Logs, Fetch Changed Files, Build AI Prompt, Get Model URL, Get File SHA, Push Fix File and Post Failure Comment). If one of them fails, n8n records an execution error but nothing is written to GitHub (source: docs/\allowbreak{}architecture/\allowbreak{}n8n-workflow.md, "Known limits"). The Round 5 notes also say that a failure after the fix is pushed, between Create Branch and Open Pull Request, is not covered by a comment. We did not re-check that sentence against later changes.

\section{Retry behaviour}

\begin{figure}[!htbp]
\centering
\begin{tikzpicture}[node distance=0.8cm and 1.2cm, every node/.style={font=\sffamily\footnotesize}]
  \node[start] (a) {CI fails on \texttt{main}};
  \node[proc, below=of a] (b) {n8n heals, attempt 1\\(thinking off)};
  \node[proc, below=of b] (c) {Fix PR opened on \texttt{ai-fix-*}};
  \node[proc, below=of c] (d) {CI runs on the PR};
  \node[decision, below=of d] (e) {CI passes?};
  \node[start, right=2.8cm of e] (f) {Done: a person\\reviews and merges};
  \node[decision, left=2.4cm of e] (g) {attempts\\left?};
  \node[proc, above=1.35cm of g] (h) {n8n heals again\\(thinking on), same branch};
  \node[start, below=of g] (i) {Comment: a person\\needs to look};
  \draw[flow] (a) -- (b); \draw[flow] (b) -- (c); \draw[flow] (c) -- (d); \draw[flow] (d) -- (e);
  \draw[flow] (e) -- node[lbl, above]{yes} (f);
  \draw[flow] (e) -- node[lbl, above]{no} (g);
  \draw[flow] (g) -- node[lbl, right]{yes} (h);
  \draw[flow] (h) |- (d);
  \draw[flow] (g) -- node[lbl, right]{no} (i);
\end{tikzpicture}
\caption{The capped retry loop: a failing heal pull request wakes the workflow again, and the attempt number is carried in the commit trailer.}
\label{fig:retry}
\end{figure}


The retry loop is driven by GitHub Actions, not by anything inside n8n. When CI fails on a heal PR, the notify job sends the PR's head branch, its head commit and the trailer number as \texttt{attempt}. In n8n this has four effects. The prompt carries the "Heal attempt N+1" line. Thinking is switched on. The file is fetched at the PR head, so it holds the previous fix. And the new commit goes onto the same branch with the next trailer number. The cap is checked only in the notify job. With \texttt{MAX\_\allowbreak{}ATTEMPTS} set to 3, the first heal and two retries make three attempts, and then the notify job comments that a human is needed.

One detail of the retry shows up at the last step. The new commit is on the PR's own branch, which is also the base that Open Pull Request is given, so GitHub answers 422 "No commits between X and X". Heal Outcome treats this as normal for attempt 2 or later and reports \texttt{retry\_\allowbreak{}pushed}. The need for this was found live: all three retries of one round had pushed correctly and were still marked as failed (source: tools/\allowbreak{}nodes/\allowbreak{}heal\_\allowbreak{}outcome.js, comment; docs/\allowbreak{}results/\allowbreak{}GOLIVE\_\allowbreak{}RESULTS.md, Round 3).

\section{Idempotent re-runs}

The same failed commit can reach the workflow twice, for example when someone re-runs the CI job. The workflow is built so that the second run updates the first one's work instead of breaking. The branch name ends in the first 7 characters of the commit, so the same commit gives the same name when the model repeats its choice. A repeat of the same commit lands on the same branch (a measured case, according to a comment in tools/\allowbreak{}nodes/\allowbreak{}extract\_\allowbreak{}ai\_\allowbreak{}response.js). Create Branch tolerates the 422, Get File SHA reads the file's current blob on that branch, Push Fix File updates it, and if the PR already exists, Heal Outcome reports \texttt{pr\_\allowbreak{}already\_\allowbreak{}open} and leaves it alone.

This holds only when the model proposes the same branch name. Because the name comes from the model, a different name for the same commit would open a second branch and a second PR. We did not test that case.
